\begin{abstract}
Dynamic model ensembling in latent coordinate sandboxes has emerged as a lightweight, training-free mechanism for continuous task adaptation during inference. However, deploying such architectures on resource-constrained edge hardware requires low-precision quantization (e.g., 8-bit integer [INT8] activations and 4-bit integer [INT4] ensembling weights) to meet strict latency, compute, and memory bandwidth budgets. Under these extreme low-precision limits, standard dynamic ensembling methods experience severe representation collapse: quantization rounding noise destroys subtle activation boundaries, causing centroids to overlap, gradients to vanish, and dynamic routing to degrade to the performance of static uniform merging. To resolve this deployment bottleneck, we propose \textbf{QA-Merge} (\textbf{Q}uantization-Aware \textbf{Merge}), a pragmatically designed, hardware-compatible suite of techniques that stabilizes dynamic ensembling in low-precision spaces. First, \emph{Quantized Centroid Calibration (QCC)} computes and calibrates task signatures directly within the quantized integer representation space, preventing centroid overlap. Second, \emph{Straight-Through Estimator (STE) Gating} optimizes parametric routers through discrete rounding boundaries to avoid gradient vanishing. Third, \emph{Error-Feedback Trajectory Stabilization (EF-Smooth)} recursively diffuses weight rounding errors. Fourth, we introduce \emph{Activation Error Feedback (AEF)} to accumulate layer-wise activation rounding errors, preventing tiny dynamic updates from being rounded to zero. Evaluated on the high-fidelity Coordinate Sandbox, QA-Merge successfully recovers full-precision ensembling gains under INT8/INT4 constraints with zero downstream accuracy loss and minimal computational overhead, proving that adaptive latent ensembling is highly robust and fully deployable on low-cost edge processors.
\end{abstract}
